\section{Experimental Setup}\label{sec:setup}
% Written in P8 (framework-independent). Sources: configs/{data,models,grid,quantization,analysis}.yaml,
% run_config.json of the full-grid runs, benchmark JSON files, data_profile.csv / split_summary.csv (Table I),
% pareto.csv (budgets); decisions K-006, K-008, K-022, K-026 ... K-031, K-034 ... K-041, K-046 ... K-062.
% No numeric literal: every number is a macro.

\subsection{Datasets}
\emph{Car-Hacking}~\cite{song2020dcnn} contains \CHAttackFiles{} attack files (DoS, fuzzy, gear spoofing and RPM
spoofing) and an attack-free recording, with frame labels supplied by the dataset. \emph{ROAD}~\cite{verma2024road}
contains attack-free (ambient) captures, fuzzing, targeted fabrication attacks and their masquerade versions. ROAD
has no frame labels; a frame is labelled injected if its identifier is the injection identifier of the capture
metadata, its payload matches the injected pattern and it falls inside the injection interval (fuzzing: an
all-ones payload after the start of the injection). The \ExcludedROAD{} accelerator captures, whose injection
happened before the recording started, have no injected frames to label and are excluded. Table~\ref{tab:dataset}
lists frames, injected frames and windows per attack type and split; attacks are rare at window level, especially
in ROAD.

\begin{table*}[t]
\centering
\footnotesize
\caption{Datasets and splits. Captures: captures that contribute windows of the row to the training, validation and test split (attack-free windows also come from the attack-free parts of attack captures). Frames and injected frames: whole captures of that attack type. Windows of N = 100 frames, stride 50 in training and 100 (non-overlapping) in validation and test. ROAD is split by source capture, so no capture, and no masquerade log derived from it, feeds two splits; Car-Hacking has one file per attack, which is cut into disjoint time blocks separated by discarded gaps.}
\label{tab:dataset}
\begin{tabular}{@{}llrrrrrrl@{}}
\toprule
Attack type & Group & Captures & Frames & Injected & \multicolumn{3}{c}{Windows} & Note \\
\cmidrule(lr){6-8}
& & tr/va/te & & frames & train & val & test & \\
\midrule
\multicolumn{9}{l}{\textit{Car-Hacking (split: disjoint time blocks of each file)}} \\
attack-free & -- & 6/5/5 & 1,879,939 & -- & 106,938 & 14,147 & 15,559 & frames: attack-free captures only \\
DoS & coarse & 1/1/1 & 2,774,703 & 587,521 & 13,112 & 2,096 & 2,206 &  \\
fuzzy & coarse & 1/1/1 & 2,947,792 & 491,847 & 18,667 & 1,723 & 1,962 &  \\
gear spoofing & coarse & 1/1/1 & 3,552,074 & 597,252 & 22,637 & 3,764 & 3,814 &  \\
RPM spoofing & coarse & 1/1/1 & 3,730,634 & 654,897 & 24,025 & 4,339 & 4,308 &  \\
\midrule
\multicolumn{9}{l}{\textit{ROAD (split: by source capture)}} \\
attack-free & -- & 14/11/16 & 24,149,519 & -- & 304,260 & 53,226 & 53,526 & frames: attack-free captures only \\
fuzzing & coarse & 1/0/2 & 94,931 & 1,061 & 173 & 0 & 68 & no validation capture \\
correlated signal & coarse & 1/1/1 & 192,748 & 5,493 & 1,042 & 533 & 315 &  \\
max speedometer & coarse & 1/1/1 & 572,842 & 11,694 & 1,230 & 783 & 1,522 &  \\
reverse light off & coarse & 1/1/1 & 308,385 & 5,480 & 337 & 591 & 607 &  \\
reverse light on & coarse & 1/1/1 & 465,234 & 8,034 & 995 & 920 & 587 &  \\
correlated signal (masq.) & stealthy & 1/1/1 & 187,258 & 5,493 & 1,000 & 512 & 302 &  \\
max speedometer (masq.) & stealthy & 1/1/1 & 561,151 & 11,694 & 1,181 & 752 & 1,461 &  \\
reverse light off (masq.) & stealthy & 1/1/1 & 302,908 & 5,480 & 323 & 568 & 583 &  \\
reverse light on (masq.) & stealthy & 1/1/1 & 457,203 & 8,034 & 955 & 883 & 564 &  \\
max coolant temp & coarse & 0/0/1 & 61,923 & 43 & 0 & 0 & 43 & test only (unseen in training) \\
max coolant temp (masq.) & stealthy & 0/0/1 & 61,881 & 43 & 0 & 0 & 43 & test only (unseen in training) \\
\bottomrule
\end{tabular}
\end{table*}


\subsection{Training}
Every model is trained on both datasets with \NSeeds{} seeds (\Seeds), \NRuns{} training runs in total. The deep
models use Adam with learning rate \LearningRate, weight decay \WeightDecay, batch size \BatchSize{} (1D-CNN) or
\AEBatchSize{} (conv-AE), at most \MaxEpochs{} epochs and early stopping with patience \Patience. The 1D-CNN keeps
the epoch with the highest validation average precision, which measures the ranking of attack windows independently
of the threshold, and is trained with a positive-class weight equal to the ratio of normal to attack training
windows (Car-Hacking \PosWeightCH, ROAD \PosWeightROAD); no resampling is used. Its training set omits the
attack-free windows of masquerade captures, which duplicate those of the fabrication source
(\CNNTrainWindowsCH{} and \CNNTrainWindowsROAD{} windows). The conv-AE is trained on attack-free training windows
with no frame within \NormalMarginS~s of an injected frame, again without masquerade captures (\AETrainWindowsCH{}
and \AETrainWindowsROAD{} windows), and stops on the reconstruction loss of the clean validation normals.

XGBoost uses \XGBEstimators{} trees of depth at most \XGBDepth, learning rate \XGBLR, the same class weight and
early stopping after \XGBEarlyStop{} rounds without improvement of the validation log-loss. Its row and column
sampling rates are \XGBSubsample, which makes the training deterministic: the \NSeeds{} seeds yield the same model,
so XGBoost has no seed variance and counts as a single observation wherever seeds are compared.

Training runs on an \GPUName{} with \TrainThreads{} CPU threads and deterministic settings (fixed seeds,
deterministic cuDNN kernels). Every run stores its configuration, library versions (Python \PythonVersion, PyTorch
\TorchVersion, ONNX \ONNXVersion, ONNX Runtime \ORTVersion, scikit-learn \SklearnVersion, XGBoost \XGBVersion) and
the git commit. The thresholds are fitted on \CalibNormalsCH{} (Car-Hacking) and \CalibNormalsROAD{} (ROAD) clean
validation normal windows. In Car-Hacking, the attack-free gaps between bursts are shorter than the
\NormalMarginS-s margin, so these windows all come from the attack-free recording.

\subsection{Evaluation and metrics}
Metrics are computed per attack type on the test split. The evaluation set of a type contains its attack windows
and the normal windows of the captures that contain it, so that false alarms are counted under the same vehicle
and driving conditions. The primary metric is the recall (true-positive rate) at the calibrated threshold: it
depends only on the attack windows and is therefore comparable across types whose attack prevalence differs, and
since every version runs at the same false-alarm rate it is comparable across models and quantization levels. The
quantization loss of type $t$ at level $k$ is
\begin{equation}
\dRecall_t(k) = \mathrm{Recall}_t(\Qfp) - \mathrm{Recall}_t(k),\qquad k \in \{\Qhalf, \Qint\},
\end{equation}
so a positive value is a loss. F1, ROC-AUC and average precision and their losses are secondary; the false-positive
rate is reported on the attack-free test captures. Pooled metrics over all types are reported as secondary
results only, and accuracy is not reported because of the class imbalance. The engine-coolant pair, unseen in
training, is reported apart. A paired unit whose float32 recall (mean over seeds) is below \FloorRecall{} on either
side is excluded from the paired analysis, because its loss is bounded by the floor; the rule depends on \Qfp{}
only.

\subsection{Statistical analysis}
Uncertainty is estimated with a hierarchical bootstrap that crosses training seeds and test windows. Each replicate
draws seeds with replacement and resamples the test windows once, with the same draw for every seed, quantization
level and model, so that differences between levels, between the two members of a pair and between models stay
paired. Windows are resampled in moving blocks, whose length is the cube root of the stratum size rounded up,
within each stratum of capture segment and label, because consecutive windows of one injection are correlated. We
use \NBoot{} replicates and \CIPct{} percentile intervals; thresholds are kept fixed in every replicate.

For the coarse--stealthy comparison, the primary statistic is the paired difference
$D = \dRecall_{\mathrm{masq}} - \dRecall_{\mathrm{fab}}$ per fabrication--masquerade unit and its mean over units
(a positive $D$ means a larger loss on the stealthy side). With few units no paired test can reach significance, so this analysis is
exploratory and reports effect sizes with intervals. The secondary group comparison on the ROAD primary types
(\GroupStealthy{} stealthy against \GroupCoarse{} coarse) uses the exact two-sided Mann--Whitney $U$ test, Cliff's
$\delta$ and the Holm correction, and we report the smallest attainable $p$-value of each test next to its result.
The latency--size--recall trade-off is analysed as a Pareto front; a candidate counts as dominated only if it is
dominated in both measurement sessions. The supervised and unsupervised models are compared per attack type over
the \ASThreeTypes{} types of both datasets on $\Delta$ROC-AUC, and on $\dRecall$ only where both models reach a
float32 recall of at least \RecallFloorAS.

The verdict rules for the hypotheses were written after the descriptive results of the full grid were available
but before any interval or test was computed; the robustness checks of the supervised--unsupervised comparison were
added after a smoke run of the analysis. Both are reported as such together with the results.

\subsection{Measurement environment}
Latency and memory are measured on a laptop host, an \CPUName{} (\PerfCores{} performance and \EffCores{}
efficiency cores, \PhysCores{} physical and \LogicalCores{} logical processors, \MemGiB~GiB RAM) under
\OSName. The process is pinned to one logical processor on each of \PinnedCores{} performance cores (skipping core
zero, which serves most interrupts), and the measurement starts only in the best-performance power mode. Inference
uses the CPU execution provider of ONNX Runtime \BenchORTVersion{} with \BenchIntraOp{} intra-operator and
\BenchInterOp{} inter-operator threads, sequential execution and batch size \BenchBatch. After at least
\BenchWarmup{} warm-up inferences and at least \BenchWarmupS~s, which are discarded, \BenchRepeats{} inferences are
timed, cycling through \BenchInputs{} fixed validation windows. Peak memory is measured in a separate process per
model; model size is the file size.

Two sessions are measured; they count as consistent if the median latencies differ by at most \MedianTol{} and the
\TailPct{}th-percentile latencies by at most \PNinetyFiveTol. A rule fixed before the measurement states that inconsistent
targets are not re-measured: both sessions are reported, latency comparisons rest on within-session ratios to
\Qfp{}, given as the range over the two sessions, a direction is stated only if both sessions agree, tail
percentiles are not interpreted and absolute latencies are read as orders of magnitude only. Latency is measured on
the seed-\BenchSeed{} models, since the graph, and with it the latency, depends on the architecture and the
quantization level, not on the trained weights. The float16 models are not timed (Section~\ref{sec:method}). The
host is an approximation of an edge device: no ECU, single-board computer or resource-limited container was
measured.

To relate latency to the bus, we define two budgets from the median message rate of the attack-free captures
(Car-Hacking \MsgRateCH, ROAD \MsgRateROAD{} messages per second): the \emph{window budget}, the time in which
\BudgetWindowMsgs{} messages arrive (\BudgetWindowCH~ms and \BudgetWindowROAD~ms), suits detection on
non-overlapping windows; the \emph{strict budget}, \BudgetStrictMsgs{} message interval (\BudgetStrictCH~ms and
\BudgetStrictROAD~ms), suits a sliding window that is re-evaluated on every frame.
